\documentclass[11pt,a4paper]{article}

% ---- packages -----------------------------------------------------------------
\usepackage[margin=2.4cm]{geometry}
\usepackage{amsmath,amssymb}
\usepackage{booktabs}
\usepackage{array}
\usepackage{tabularx}
\usepackage{xcolor}
\usepackage{graphicx}
\usepackage{enumitem}
\usepackage{microtype}

% colours (define BEFORE hyperref so its options resolve)
\definecolor{NavyBlue}{RGB}{30,80,150}
\definecolor{RuleGrey}{RGB}{120,120,120}
\definecolor{BoxBack}{RGB}{244,247,251}
\definecolor{BoxLine}{RGB}{200,212,228}

\usepackage[hidelinks,colorlinks=true,linkcolor=NavyBlue,urlcolor=NavyBlue,citecolor=NavyBlue]{hyperref}

% section styling
\usepackage{sectsty}
\sectionfont{\color{NavyBlue}\large}
\subsectionfont{\color{NavyBlue}\normalsize}

% headers/footers
\usepackage{fancyhdr}
\pagestyle{fancy}
\fancyhf{}
\renewcommand{\headrulewidth}{0.3pt}
\fancyhead[L]{\footnotesize\color{RuleGrey} Conflict $\to$ Road-Blockage Framework}
\fancyhead[R]{\footnotesize\color{RuleGrey} ACLED Afghanistan 2017--2026}
\fancyfoot[C]{\thepage}

% a light "insight/note" box
\usepackage[most]{tcolorbox}
\newtcolorbox{notebox}{colback=BoxBack,colframe=BoxLine,boxrule=0.6pt,
  arc=2pt,left=6pt,right=6pt,top=5pt,bottom=5pt}

% code font for inline identifiers
\newcommand{\code}[1]{\texttt{\small #1}}

% math shortcuts
\newcommand{\Reff}{R_{\mathrm{eff}}}
\newcommand{\Rphys}{R_{\mathrm{phys}}}
\newcommand{\Rgeo}{R_{\mathrm{geo}}}
\newcommand{\Msev}{M_{\mathrm{sev}}}
\newcommand{\Pblock}{P_{\mathrm{block}}}

\title{\vspace{-1.2cm}\textbf{\color{NavyBlue} A Framework for Estimating Road-Blockage
Probability from Conflict Events}\\[2pt]
\large A weakly-supervised, parametric impact model}
\author{Conflict-event $\to$ Road-Blockage (\textsc{c2rb}) model}
\date{ACLED, Afghanistan, 2017--2026 \quad(69{,}655 geocoded events)}

\begin{document}
\maketitle
\thispagestyle{fancy}
\vspace{-0.4cm}

% =====================================================================================
\section{Problem and approach}

\textbf{Goal.} Given a conflict event (or a set of events over a time window), estimate which
nearby roads are \emph{affected} and the \emph{probability that each road is blocked}. The
central modelling question is: \textbf{how large is an event's impact ``buffer zone'', and how
does it differ by \code{sub\_event\_type}?}

\textbf{The core difficulty --- no ground truth.} The ACLED export has no column that says
``a road was blocked''. A fully supervised classifier therefore cannot be trained directly.
Two naive alternatives both fail:
\begin{itemize}[nosep,leftmargin=1.4em]
  \item \emph{Pure assumption} (pick buffer radii by hand): not defensible, not reproducible.
  \item \emph{Pure ML} (train a closure classifier): impossible without labels.
\end{itemize}

\textbf{Our approach --- a weakly-supervised parametric probability model.} We build an
explicit, interpretable probability model whose parameters are (a)~\textbf{seeded from
published literature} and (b)~\textbf{calibrated against weak labels mined from the ACLED
\code{notes} free text}. This is defensible, reproducible, and mirrors established practice:
ACLED's own \emph{Conflict Exposure} product uses fixed event-type buffers (1/2/5\,km) with
Voronoi de-duplication, and a published conflict-impact model assigns per-event-type radii
(Battles 25\,km, Explosions 10\,km, Protests/Riots 5\,km) with type-specific decay shapes.

\begin{notebox}
\textbf{Pipeline.}\quad
$\text{event}\;\longrightarrow\;\Reff\;\longrightarrow\;K(d)\;\longrightarrow\;
\Pblock(d)\;\xrightarrow{\;\text{noisy-OR}\;+\;\text{time decay}\;}\;P(\text{road blocked})$
\end{notebox}

% =====================================================================================
\section{The model}

For an event $e$ with sub-event-type $s$, geo-precision code $g\in\{1,2,3\}$, fatalities
$n_f$, and civilian-targeting flag $c\in\{0,1\}$:

\subsection{Effective buffer radius --- ``how big is the zone?''}

\begin{equation}
\Reff(e) \;=\; \sqrt{\,\Rphys(s)^2 + \Rgeo(g)^2\,}\;\cdot\;\Msev(e)
\label{eq:reff}
\end{equation}

\begin{itemize}[nosep,leftmargin=1.4em]
  \item $\Rphys(s)$ --- physical/operational impact radius of the event type (km),
        literature-seeded per \code{sub\_event\_type}. E.g.\ IED 1\,km, armed clash 15\,km,
        protest 4\,km.
  \item $\Rgeo(g)$ --- \emph{locational uncertainty} of where the event actually happened.
        ACLED geo-precision $1\approx$ town (1\,km), $2\approx$ near a town / part of a region
        (5\,km), $3\approx$ wider region (25\,km). A low-precision event must spread its
        blockage probability over a larger uncertain area.
  \item \textbf{Quadrature combination} $\sqrt{\Rphys^2+\Rgeo^2}$. Both terms behave like
        independent Gaussian length-scales (physical spread and locational error); variances of
        independent Gaussians add, so their standard deviations add in quadrature. This is more
        principled than naive addition (which double-counts) and never smaller than either term.
\end{itemize}

The severity multiplier $\Msev(e)$ stretches the radius for deadlier and civilian-directed
events, with diminishing returns so a 50-fatality battle does not produce a $50\times$ radius:
\begin{align}
S(n_f)   &= \alpha\,(1+n_f) + (1-\alpha)\bigl(1+\ln(1+n_f)\bigr) \label{eq:sev}\\[2pt]
\Msev(e) &= \operatorname{clip}\!\Bigl(1 + \kappa\Bigl(\tfrac{S(n_f)}{S(n_{\mathrm{ref}})}-1\Bigr),\,
            0.5,\,3.0\Bigr)\,\cdot\,(1+\gamma_{\mathrm{civ}}\,c) \label{eq:msev}
\end{align}
with defaults $\alpha=0.15$ (mostly log growth), $\kappa=0.25$, $\gamma_{\mathrm{civ}}=0.30$.
The \code{clip} bounds keep outliers from exploding the buffer.

\subsection{Distance-decay blockage probability --- ``how does it fall with distance?''}

\begin{equation}
\Pblock(d,e) \;=\; P_0(s)\,\cdot\,K_s\!\left(d;\,\Reff\right)
\label{eq:pblock}
\end{equation}

\begin{itemize}[nosep,leftmargin=1.4em]
  \item $P_0(s)$ --- \emph{peak blockage propensity}: the probability that an event of type $s$
        blocks a road at the epicentre. \textbf{This is the parameter calibrated from
        \code{notes}} (Sec.~\ref{sec:calib}).
  \item $K_s(d)\in[0,1]$ --- distance-decay kernel, its \emph{shape} chosen by the event
        type's physics (Table~\ref{tab:kernels}).
\end{itemize}

\begin{table}[h]
\centering
\renewcommand{\arraystretch}{1.35}
\begin{tabularx}{\textwidth}{@{}l l l X@{}}
\toprule
\textbf{Kernel} & \textbf{Formula} & \textbf{Used for} & \textbf{Rationale} \\
\midrule
Gaussian & $\exp\!\left(-\dfrac{d^2}{2\sigma^2}\right),\ \sigma=\Reff/2$
  & IED, suicide, grenade, shelling, airstrike & sharp, blast-like, fast decay \\
Exponential & $\exp\!\left(-\dfrac{d}{\lambda}\right),\ \lambda=\Reff/3$
  & armed clash / battles & heavier tail --- fighting spreads along terrain/roads \\
Uniform & $\mathbf{1}\{d\le \Reff\}$
  & protests, riots, territorial control & roughly constant over an area, then cut off \\
Point & $\mathbf{1}\{d\le \max(\Reff,0.3)\}$
  & targeted violence, arrests & effect confined to the site \\
\bottomrule
\end{tabularx}
\caption{Distance-decay kernels, chosen per event type.}
\label{tab:kernels}
\end{table}

\subsection{Combining many events --- noisy-OR with temporal decay}

A road segment near several events is blocked if \emph{any} of them blocks it. Treating events
as independent blocking attempts gives the \textbf{noisy-OR}, with an \textbf{exponential
temporal decay} $\varphi(\Delta t)=0.5^{\,\Delta t/\tau}$ (half-life $\tau=30$ days) so older
events fade:
\begin{equation}
P_{\text{segment}} \;=\; 1 - \prod_{e\,\in\,W}\Bigl[\,1 - \Pblock(d_e,e)\,\cdot\,\varphi(\Delta t_e)\,\Bigr]
\label{eq:noisyor}
\end{equation}
over a time window $W$. ACLED-style Voronoi de-duplication can be layered on to avoid
double-counting tightly-clustered events (noted as an extension).

\subsection{From points to roads}

We use the WFP/AGCHO Afghanistan road network
(\code{afg\_trs\_roads\_l\_arazi\_2023x\_fixed\_segmented.shp}): \textbf{27{,}991 LineString
segments, $\approx$1\,km each} (median 1{,}068\,m), EPSG:4326, each carrying road name, class,
average slope and altitude. Each segment is represented by its midpoint; $d_e$ is the
segment-to-event distance. Every segment within $r_{\max}=3\,\Reff$ of an event receives that
event's $\Pblock(d_e,e)$, and segments are aggregated by noisy-OR. The result is a per-segment
blockage probability that rolls up to \textbf{named roads}. Without a road layer the same
scoring runs over a regular lat/lon grid, producing a continuous blockage-risk raster (the
fallback) --- the code path is identical, only the target points differ.

% =====================================================================================
\section{Calibration \& validation from \texttt{notes} (weak supervision)}
\label{sec:calib}

\subsection{Building a weak ``road-disruption'' label}

The export has no ``road blocked'' column, so we create one from the \code{notes} free text.
Reading the notes reveals \textbf{three distinct mechanisms} by which an event disrupts a road;
each is captured by its own regular expression and the label is their union,
$\text{road\_disruption}=\text{block}\lor\text{hazard}\lor\text{denial}$.

\begin{table}[h]
\centering
\renewcommand{\arraystretch}{1.25}
\begin{tabularx}{\textwidth}{@{}l X r l@{}}
\toprule
\textbf{Mechanism} & \textbf{Signal in notes} & \textbf{Coverage} & \textbf{Typical type} \\
\midrule
Blockage & ``blocked/closed the highway'', ``cut off the road'', ``under siege'' & 0.14\% & protests, sieges \\
Hazard   & ``roadside bomb'', ``IED/landmine on the road'', ``bridge destroyed'' & 2.36\% & IED / landmine \\
Denial   & ``convoy ambushed'', ``set up checkpoint'', ``seized control of the road'' & 1.89\% & armed clashes \\
\midrule
\textbf{Union} & any of the above & \textbf{4.35\%} & 3{,}032 events \\
\bottomrule
\end{tabularx}
\caption{The three road-disruption mechanisms mined from \code{notes}.}
\label{tab:mechanisms}
\end{table}

Crucially the label is derived \textbf{only from text, never from \code{sub\_event\_type}}
(only $\sim$18\% of IEDs are ``roadside''), so the per-type rates below are genuine signal, not
a relabelling of the type. This three-mechanism distinction is itself a finding: \emph{road
hazard} (a mined road) and \emph{deliberate blockage} (protesters on a highway) are different
phenomena that a single ``road blocked?'' keyword would have missed.

\subsection{Peak blockage propensity $P_0(s)$}

$P_0(s)=P(\text{road\_disruption}\mid s)$ is estimated with \textbf{Beta--Binomial
empirical-Bayes shrinkage} (prior strength 50, anchored at the 4.35\% global rate) so
small-sample types (Grenade $n{=}176$, Suicide $n{=}289$) get stable estimates instead of
noisy zeros:
\begin{equation}
\hat{P}_0(s) \;=\; \frac{k_s + a_0}{n_s + a_0 + b_0},
\qquad a_0 = \bar p\,m,\quad b_0=(1-\bar p)\,m,\quad m=50,
\end{equation}
where $k_s$ positives out of $n_s$ events of type $s$, and $\bar p$ is the global rate.
Selected results with Wilson 95\% intervals (full table in \code{params.yaml}):

\begin{table}[h]
\centering
\renewcommand{\arraystretch}{1.2}
\begin{tabular}{@{}l r r l@{}}
\toprule
\textbf{sub\_event\_type} & \textbf{$n$} & \textbf{$P_0$ (shrunk)} & \textbf{95\% CI} \\
\midrule
Remote explosive/landmine/IED      & 8{,}545  & \textbf{0.181} & [0.174, 0.190] \\
Suicide bomb                       & 289      & 0.098 & [0.077, 0.148] \\
Armed clash                        & 39{,}974 & 0.031 & [0.029, 0.032] \\
Peaceful protest                   & 1{,}454  & 0.019 & [0.012, 0.026] \\
Shelling/artillery/missile attack  & 2{,}540  & 0.010 & [0.007, 0.014] \\
Air/drone strike                   & 6{,}249  & \textbf{0.002} & [0.001, 0.003] \\
\bottomrule
\end{tabular}
\caption{Calibrated blockage propensity --- the headline result.}
\label{tab:p0}
\end{table}

\begin{notebox}
\textbf{Headline result.} Blockage propensity spans $\sim$90$\times$ across types: IEDs and
landmines --- deliberately placed \emph{on roads} to deny movement --- dominate, while
air/drone strikes (which target buildings and people) almost never produce road-disruption
language. This empirically confirms that the buffer/blockage model \textbf{must} vary by
\code{sub\_event\_type}.
\end{notebox}

\subsection{Validation --- does the signal generalise?}

A logistic regression $\text{road\_disruption}\sim\code{sub\_event\_type}+\log(1+\text{fatalities})
+\code{geo\_precision}+\code{civilian\_targeting}$, evaluated with 5-fold cross-validation:

\begin{itemize}[nosep,leftmargin=1.4em]
  \item \textbf{AUC $= 0.79$} --- event attributes alone discriminate road-disrupting events
        well above chance.
  \item \textbf{Brier $= 0.039$} --- \emph{better} than the no-skill baseline (0.042), i.e.\
        genuinely informative.
  \item \textbf{Well-calibrated} --- predicted $\approx$ observed across all deciles (top decile
        predicts 0.19, observes 0.17). Calibration requires the plain model;
        \code{class\_weight="balanced"} improves ranking but inflates probabilities (use
        isotonic/Platt calibration if class-weighting).
\end{itemize}

The model is both discriminative and calibrated, so the $P_0(s)$ values are trustworthy as
probabilities, not merely rankings.

\subsection{Other sanity checks (in the notebook)}
\begin{itemize}[nosep,leftmargin=1.4em]
  \item \textbf{Roads-within-1\,km benchmark --- passed.} With the real road network loaded,
        \textbf{78\% of events fall within 1\,km of a road}, independently matching the
        published $\sim$70\% figure --- strong external validation.
  \item \textbf{Sensitivity analysis} over $\Rphys$ and decay-kernel choice shows how the
        affected-road footprint responds to assumptions (what is assumed vs.\ learned).
  \item \textbf{Severity \& geo-precision} move the weak label in sensible directions.
\end{itemize}

% =====================================================================================
\section{Statistical methods used}

Distance-decay kernels (Gaussian / exponential / uniform); quadrature combination of
independent uncertainty scales; bounded log-linear severity scaling; noisy-OR probabilistic
aggregation with exponential temporal decay; weak supervision via regex labelling of free text;
Beta--Binomial empirical-Bayes shrinkage with Wilson confidence intervals; logistic regression
with cross-validated AUC / Brier / calibration-curve diagnostics; (optional) Voronoi
tessellation for overlap de-duplication and KDE for hotspot context.

% =====================================================================================
\section{Limitations \& honest caveats}

\begin{enumerate}[leftmargin=1.6em]
  \item \textbf{Weak labels are a proxy, not ground truth.} \code{road\_disruption} measures
        \emph{reported} road-relevant language in \code{notes}, not verified physical closures.
        Reporting bias propagates into $P_0$. Replace with real closure data
        (OCHA/Logistics Cluster access reports) when available --- the calibration code is
        unchanged.
  \item \textbf{$\Rgeo$ cannot be learned from this data.} ACLED snaps all events at a named
        \code{location} to one shared centroid, so intra-location coordinate spread is
        structurally $\sim$0\,km. We therefore fall back to ACLED's \emph{documented} precision
        semantics (1/5/25\,km). With finer coordinates, the empirical-spread estimator in the
        notebook would populate $\Rgeo$.
  \item \textbf{$\Rphys$ and kernel shapes are assumptions} (literature-seeded), unlike $P_0$
        (data-calibrated). The sensitivity analysis quantifies how much they matter; tune in
        \code{params.yaml}.
  \item \textbf{Independence in noisy-OR} slightly over-counts spatially correlated events;
        Voronoi de-duplication (ACLED's approach) is the recommended refinement.
  \item \textbf{Static / impact-only.} This scores the impact of given events; it does not
        forecast where future events will occur.
\end{enumerate}

% =====================================================================================
\section{Files}

\begin{table}[h]
\centering
\renewcommand{\arraystretch}{1.2}
\begin{tabularx}{\textwidth}{@{}l X@{}}
\toprule
\textbf{File} & \textbf{Purpose} \\
\midrule
\code{c2rb.py} & Framework core: weak-label regexes, calibration, radius/kernel/noisy-OR math, scorer \\
\code{params.yaml} & All tunable parameters: $\Rphys$, decay class, $\Rgeo$, calibrated $P_0$, severity \& time constants \\
\code{conflict\_road\_blockage.ipynb} & Annotated end-to-end walkthrough: load $\to$ label $\to$ calibrate $\to$ score $\to$ map $\to$ sensitivity \\
\code{build\_notebook.py} & Regenerates the notebook from cell definitions (diffable source of truth) \\
\code{requirements.txt} & Pinned dependencies for reproducibility \\
\code{methodology.tex / .pdf} & This document \\
\bottomrule
\end{tabularx}
\end{table}

\textbf{Reproduce.} Create the venv and \code{pip install -r requirements.txt}, then run the
notebook top-to-bottom. The road and admin layers are already wired in \code{params.yaml}; set
\code{paths.roads: null} to fall back to grid-raster mode.

% =====================================================================================
\section{References}

\begin{itemize}[nosep,leftmargin=1.4em]
  \item ACLED, \emph{Conflict Exposure methodology} (buffer zones 1/2/5\,km, event-type
        adjustment, Voronoi de-duplication).\\
        \url{https://acleddata.com/methodology/conflict-exposure-methodology}
  \item ACLED, \emph{Codebook} (event types, sub-event types, geo-precision codes) ---
        \code{AcledCodebook.md}.
  \item \emph{Applying Computational Engineering Modelling to Analyse the Social Impact of
        Conflict and Violent Events} (per-type impact radii \& decay shapes), PMC12562815.\\
        \url{https://pmc.ncbi.nlm.nih.gov/articles/PMC12562815/}
  \item \emph{Proximity to roads shapes patterns of violence in North and West Africa}
        ($\sim$70\% of violent events within 1\,km of a road).\\
        \url{https://anl.geog.ufl.edu/roads-violence-nwafrica/}
  \item Road network: WFP / AGCHO-NSIA Afghanistan segmented roads
        (\code{afg\_trs\_roads\_l\_arazi\_2023x}); admin boundaries: AGCHO-NSIA Afghanistan
        adm1/adm2 (2017).
\end{itemize}

\end{document}
